\chapter{La valutazione dei sistemi energetici}
\label{cap:sistemienergetici}

Argomento conclusivo del modulo, collocato per ultimo per un duplice motivo:
\begin{itemize}
  \item il \textbf{sistema metabolico è conseguenza dell'attivazione del sistema neuromuscolare}:
    è il tipo di unità motorie reclutate --- e quindi l'intensità --- a determinare quale sistema
    energetico venga utilizzato;
  \item l'argomento verrà poi ripreso in chiave applicativa nelle parti successive del corso.
\end{itemize}
La fibra muscolare ricostituisce l'ATP per tre vie metaboliche, distinte per \textbf{potenza} e per
\textbf{durata} del contributo; a ciascuna corrispondono test di valutazione propri, di laboratorio
o da campo.

% ---------------------------------------------------------------------------
\section{Il metabolismo energetico della fibra muscolare}

\rubrica{Lo schema d'insieme}
Lo schema raccoglie tutti i sistemi in grado di produrre l'ATP utile alla prestazione --- cioè allo
scorrimento dell'actina sulla miosina e alla formazione dei ponti trasversali --- e distingue gli
scambi con il capillare dai due comparti interni. Lo \textbf{spessore delle frecce} indica la
potenza della rispettiva via metabolica:
\begin{enumerate}
  \item \textbf{prima fonte}: l'ATP già presente, scisso in ADP + P dall'\textit{ATPasi miosinica}
    dell'elemento contrattile;
  \item \textbf{pool dei fosfati}: il \textbf{creatinfosfato} cede lo ione fosfato all'ADP; la
    \textbf{creatina} torna poi al mitocondrio a riprendere un altro ione fosfato, e il ciclo
    ricomincia;
  \item \textbf{sistema glicolitico}: interviene quando i livelli di ADP diventano eccessivi, cioè
    quando il pool dei fosfati non è più in grado di sostenere il gradiente di forza richiesto;
  \item \textbf{comparto aerobico} (mitocondrio): acidi grassi liberi e acidi grassi in goccioline
    alimentano il \textbf{ciclo dell'acido citrico}; la \textbf{catena respiratoria} rigenera ATP da
    ADP consumando $O_2$ e liberando $CO_2$ e $H_2O$.
\end{enumerate}
Gli scambi con il capillare comprendono, in ingresso, glucosio e acidi grassi più l'$O_2$ ceduto dai
globuli rossi; in uscita lattato, ammoniaca ($NH_3$), $CO_2$ e $H_2O$.

\rubrica{Il contributo dei tre sistemi nel gesto massimale}
Nello studio di Hirvonen e coll.\ (1987) si esegue una \textbf{biopsia ogni 20\,m} su centometristi
e si misurano, in funzione della distanza percorsa alla massima velocità, i pool fosforici
intracellulari (ATP + CP), il lattato ematico e il pH:
\begin{itemize}
  \item l'\textbf{ATP} resta pressoché costante per tutta la durata della prova;
  \item il \textbf{creatinfosfato} è utilizzato moltissimo già nel riscaldamento e nella fase
    iniziale, poi cala in modo netto;
  \item il \textbf{lattato} cresce e diventa il sistema energetico prevalente: è già elevato dopo i
    primi 20--30\,m e supera gli \textbf{8\,mmol/l} a fine prova;
  \item il \textbf{pH} scende parallelamente all'aumento del lattato (7,40 $\rightarrow$ 7,20 circa);
  \item la \textbf{velocità} cresce nella prima parte, raggiunge il massimo e poi decresce.
\end{itemize}
Conclusione: uno sprint di 100\,m è \textbf{prevalentemente anaerobico lattacido}. Il dato serve
però soprattutto a capire che cosa accade nei \textbf{primi 20--30\,m}, che sono le distanze
tipiche del calciatore --- il quale usa anzi anche percorrenze più brevi, con cambi di direzione.

\rubrica{Potenza e durata dei tre sistemi}
\begin{itemize}
  \item \textbf{ATP-CP}: durata brevissima, potenza elevatissima;
  \item \textbf{anaerobico lattacido}: durata maggiore, potenza ancora elevata;
  \item \textbf{aerobico}: durata teoricamente infinita, potenza di gran lunga inferiore.
\end{itemize}
Il continuum è illustrato associando a ciascun sistema una prova dell'atletica: gesto esplosivo di
salto, corsa a ostacoli, maratona.

% ---------------------------------------------------------------------------
\section{Il sistema anaerobico alattacido (ATP-CP)}

\textbf{Meccanismo}: nel citosol il pool di \textbf{fosfocreatina} (PCr) rifosforila l'ADP prodotto
dall'idrolisi dell'ATP: $ADP + PCr \rightarrow ATP + Cr$. È la via più rapida e la prima a essere
chiamata in causa.

\rubrica{Il test di Margaria (1966)}
Messo a punto dal fisiologo italiano nel 1966, valuta la potenza anaerobica alattacida facendo
salire all'atleta \textbf{9 scalini} il più velocemente possibile:
\begin{itemize}
  \item \textbf{rincorsa}: $d$ = 6\,m dal primo scalino;
  \item \textbf{cronometraggio}: il tempo è registrato \textbf{dal terzo al nono scalino}, non
    sull'intera salita;
  \item \textbf{dislivello}: $H = \Sigma_{3-9}\,h_{sc}$, la distanza verticale data dalla sommatoria
    delle altezze dei singoli scalini;
  \item \textbf{potenza}: $P(W) = (F \times H)\,/\,T$, cioè il lavoro effettuato in funzione del
    tempo;
  \item \textbf{conversione}: la potenza veniva misurata in kg $\times$ m $\times$ min$^{-1}$ e oggi
    in watt, con 1\,W = 6,12\,kg $\times$ m $\times$ min$^{-1}$. L'equazione è riportata proprio per
    poter \textbf{confrontare dati espressi in unità di misura differenti}.
\end{itemize}

% ---------------------------------------------------------------------------
\section{Il sistema glicolitico (anaerobico lattacido)}

Via citosolica che parte dai carboidrati; il destino dell'acido piruvico dipende dall'intensità:
\begin{itemize}
  \item \textbf{substrati e resa}: dal \textbf{glicogeno muscolare} a glucosio-6-fosfato e poi a
    piruvato, con liberazione di ATP (\textbf{glucosio} $\rightarrow$ 2 ATP; \textbf{glicogeno}
    $\rightarrow$ 3 ATP);
  \item \textbf{intensità elevata}: il piruvato è trasformato in \textbf{acido lattico}, che lascia
    la cellula e si ritrova nel capillare;
  \item \textbf{intensità medio-alta}, gestibile dal sistema aerobico: il piruvato è trasformato in
    \textbf{acetil-CoA} ed entra nel \textbf{ciclo di Krebs}, molto più redditizio (34 ATP) ma a
    scapito della potenza muscolare.
\end{itemize}

\rubrica{Sistemi energetici e tipo di unità motoria}
Il pool dei fosfati e la glicolisi anaerobica --- cioè le vie con produzione di lattato --- sono i
sistemi utilizzati soprattutto dalle \textbf{fibre veloci}, quindi dalle \textbf{unità motorie
veloci}. È l'aggancio tra reclutamento e metabolismo: il grafico dell'intensità dello stimolo e
dell'attivazione muscolare è lo stesso già discusso nel capitolo \ref{cap:neuromuscolare}.

\rubrica{La riconversione del lattato e le sue implicazioni per l'allenamento}
\begin{itemize}
  \item nella prestazione degli sport di squadra l'aspetto fondamentale è la capacità di esprimere
    \textbf{altissima potenza} in condizioni anaerobiche, quindi producendo acido lattico;
  \item l'acido lattico può però essere \textbf{ritrasformato in acido piruvico} e riassorbito dal
    sistema muscolare;
  \item la conversione nelle due direzioni dipende dagli \textbf{isoenzimi della lattato
    deidrogenasi} --- le forme \textbf{1, 2, 3} da un lato e le forme \textbf{4, 5} dall'altro ---
    e l'allenamento specifico produce esattamente un \textbf{adattamento di tipo enzimatico};
  \item da qui la scelta dei \textbf{sistemi intervallati}: fasi estremamente attive di durata
    variabile alternate a fasi di recupero passivo, per creare l'adattamento biochimico;
  \item la \textbf{corsa lenta e prolungata} non trova invece impiego negli sport di squadra se non
    in defaticamento, e non come mezzo condizionante: allena solo la via che in partita non viene
    usata, perché in partita l'obiettivo è arrivare sulla palla il più presto possibile, cioè
    utilizzare il sistema anaerobico lattacido.
\end{itemize}

\rubrica{Il fattore limitante nel calciatore}
Riprendendo i fattori limitanti del capitolo \ref{cap:neuromuscolare}:
\begin{itemize}
  \item la prestazione comprende anche molta corsa a bassa intensità, ma il \textbf{fattore
    limitante} è il condizionamento delle \textbf{unità motorie veloci}, sul piano neuromuscolare
    \textbf{e} su quello metabolico;
  \item le fibre lente e intermedie partecipano invece alla \textbf{fase di recupero}, potendo
    utilizzare il sistema aerobico pur producendo una potenza relativa;
  \item tutti i \textbf{lavori submassimali} sono quindi funzionali alle fasi di recupero: saper
    alternare e condizionare il sistema neuromuscolare in modo specifico rispetto alle richieste
    della gara va letto attraverso il reclutamento delle unità motorie.
\end{itemize}

% ---------------------------------------------------------------------------
\section{Il sistema aerobico}

\rubrica{La potenza aerobica}
\begin{itemize}
  \item \textbf{valore di riferimento}: un calciatore ha normalmente una potenza aerobica di
    \textbf{50--55\,ml/kg/min};
  \item le prestazioni che richiedono una \textbf{produzione di energia prolungata} devono usare la
    via aerobica per produrre ATP;
  \item \textbf{trasporto dell'ossigeno}: l'$O_2$ si lega all'emoglobina nel polmone, è pompato dal
    cuore ai muscoli e raggiunge il mitocondrio; l'anidride carbonica torna per via venosa ed è
    liberata nel polmone;
  \item \textbf{sede}: i \textbf{mitocondri}, strutture cellulari altamente specializzate con enzimi
    ossidativi specifici; vi si metabolizzano non solo il piruvato --- e quindi il glucosio --- ma
    anche \textbf{proteine} e \textbf{grassi};
  \item \textbf{processo}: con ossigeno sufficientemente disponibile il piruvato è convertito in
    \textbf{acetil-CoA}, che entra nel ciclo di Krebs; il \textbf{sistema di trasporto degli
    elettroni} produce una grande quantità di ATP, con $CO_2$ e $H_2O$ come prodotti di scarto
    rimovibili.
\end{itemize}

\rubrica{I metabolimetri}
Lo strumento con cui si valuta la potenza aerobica analizza i \textbf{gas espirati}:
\begin{itemize}
  \item il primo dispositivo è il \textbf{sacco di Douglas}, strumentazione pionieristica e di
    difficile applicazione sul campo: si lavorava in laboratorio, dove \textbf{la prestazione era
    adeguata allo strumento di valutazione e non viceversa} --- l'inverso di come dovrebbe essere;
  \item la tecnologia ha poi reso i metabolimetri \textbf{portatili}, oggi di peso ridottissimo;
  \item resta però l'inconveniente della \textbf{maschera sul viso}, necessaria a captare i gas
    espirati: in molte discipline è di fatto impossibile usarli durante la gara o l'allenamento,
    diversamente dai cardiofrequenzimetri o dal GPS.
\end{itemize}
Il consumo di ossigeno si ricava dalla differenza tra ossigeno inspirato ed espirato:
$$V'O_2 = V'_i \cdot F_iO_2 - V'_e \cdot F_eO_2$$

\rubrica{Stato stazionario, deficit di ossigeno e debito}
\begin{enumerate}
  \item \textbf{obiettivo}: lo \textit{steady state}, condizione in cui il fabbisogno energetico è
    soddisfatto per via aerobica;
  \item \textbf{avvio dell'esercizio} (primi 2--4 minuti, riscaldamento compreso): le risposte
    fisiologiche aumentano la domanda metabolica di ossigeno, ma l'energia è fornita dal
    metabolismo anaerobico, più rapido $\rightarrow$ il pool dei fosfati si esaurisce subito,
    interviene l'acido lattico e si crea il \textbf{deficit di ossigeno};
  \item \textbf{regime}: se l'intensità non è troppo alta rispetto alla capacità di fornire ossigeno
    al muscolo si raggiunge lo stato stazionario --- è il momento in cui, in termini pratici, ``si
    rompe il fiato''. Il sistema aerobico si attiva \textbf{sempre in seguito} all'attivazione del
    sistema anaerobico lattacido;
  \item \textbf{recupero}: al termine dell'esercizio le richieste di ossigeno tornano ai valori di
    riposo, ma il consumo scende \textbf{lentamente}, perché va \textbf{ripagato il debito}
    contratto nella fase iniziale: l'energia serve a ricostituire il pool fosforico, eliminare il
    lattato accumulato e ripristinare l'omeostasi.
\end{enumerate}
Nell'esempio grafico il test è condotto a un'intensità pari al \textbf{50\,\% del $VO_2max$}, con
inizio intorno al terzo minuto e fine intorno al sedicesimo.

\rubrica{Il massimo consumo di ossigeno ($VO_2max$)}
\begin{itemize}
  \item \textbf{condizione}: quando l'intensità è così alta che le richieste metaboliche non possono
    essere soddisfatte in condizioni aerobiche --- non si arriva allo stato stazionario --- il
    muscolo integra la produzione di ATP con il metabolismo anaerobico, e la prestazione finisce
    rapidamente;
  \item \textbf{definizione}: le richieste metaboliche massime definiscono il $VO_2max$, o
    \textbf{massima potenza aerobica}: il valore più elevato di ossigeno per unità di tempo che il
    corpo umano è in grado di respirare introducendo aria;
  \item \textbf{determinanti}:
    \begin{enumerate}
      \item la \textbf{capacità metabolica cardio-polmonare}, dimensionata principalmente dalla
        \textbf{gittata cardiaca} (e dalla gittata sistolica);
      \item l'\textbf{utilizzo del sistema periferico}, cioè la differenza artero-venosa di
        pressione parziale di $O_2$;
      \item la \textbf{capacità di lavoro aerobico} del muscolo scheletrico;
    \end{enumerate}
  \item \textbf{determinazione}: incrementando progressivamente il lavoro, il $VO_2$ cresce in modo
    quasi lineare fino a stabilizzarsi su un \textbf{plateau}; il valore del plateau è il $VO_2max$
    (nell'esempio, circa 40\,ml/min/kg).
\end{itemize}

\rubrica{La velocità massimale aerobica}
In un lavoro a rampa con intensità crescente si utilizza il sistema aerobico, ma aumentando ancora
il carico --- in velocità o in watt --- si arriva a un punto in cui vengono reclutate, in sequenza,
fibre lente, fibre intermedie e infine \textbf{unità motorie veloci}, con produzione di acido
lattico. La velocità corrispondente è la \textbf{velocità massimale aerobica} (VAM).

% ---------------------------------------------------------------------------
\section{I protocolli dei test incrementali}

\rubrica{Modalità di incremento del carico (test triangolare)}
\begin{itemize}
  \item \textbf{rampa}: aumento del carico continuo e progressivo;
  \item \textbf{step}, con incremento a gradini di ampiezza proporzionale alla durata:
    \begin{itemize}
      \item \textbf{step 1 min}: +20\,W oppure +1\,km/h;
      \item \textbf{step 2 min}: +40\,W oppure +2\,km/h;
      \item \textbf{step 3 min}: +60\,W oppure +3\,km/h.
    \end{itemize}
\end{itemize}

\rubrica{Protocolli su tappeto (variazione dell'inclinazione)}
Si può incrementare l'inclinazione oppure la velocità, purché in maniera costante e fino a
esaurimento:
\begin{itemize}
  \item \textbf{test di Astrand}: velocità costante di 5\,mp/h; dopo 3$'$ in piano, l'inclinazione
    aumenta del 2,5\,\% ogni 2,5 minuti;
  \item \textbf{test di Bruce}: sia la velocità sia l'inclinazione vengono variate ogni 3 minuti;
  \item \textbf{test di Balke}: dopo 1$'$ in piano e 1$'$ a un'inclinazione del 2\,\%,
    l'inclinazione viene incrementata dell'1\,\% ogni minuto mantenendo la velocità costante a
    3,3\,mp/h.
\end{itemize}

% ---------------------------------------------------------------------------
\section{I test da campo per la determinazione indiretta del \texorpdfstring{$VO_2max$}{VO2max}}

\rubrica{Test di Cooper}
\begin{itemize}
  \item \textbf{sede e modalità}: pista di atletica leggera, corsa alla massima velocità;
  \item \textbf{durata}: 12 minuti; \textbf{parametro misurato}: distanza massima percorsa;
  \item \textbf{stima} (distanza $d$ in km):
    \begin{itemize}
      \item non allenati: $VO_2max = 22{,}351 \cdot d - 11{,}288$;
      \item allenati: $VO_2max = 11 \cdot d + 21{,}9$;
    \end{itemize}
  \item \textbf{limite}: è stato uno dei primi test usati nel calcio, ma si tratta di corsa lenta,
    o comunque a \textbf{velocità costante} da mantenere per 12 minuti --- richiesta difficile per i
    giocatori di sport di squadra e, soprattutto, \textbf{uno sforzo che in partita non si verifica
    mai}.
\end{itemize}

\begin{table}[htbp]
  \centering
  \small
  \renewcommand{\arraystretch}{1.3}
  \begin{tabularx}{\textwidth}{|c|X|c|X|}
    \hline
    \multicolumn{2}{|c|}{\textbf{Scala generale}} &
    \multicolumn{2}{c|}{\textbf{Atleti} (Bosco, 1990)} \\
    \hline
    \textbf{Metri percorsi} & \textbf{Giudizio} & \textbf{Metri percorsi} & \textbf{Giudizio} \\
    \hline
    $<$ 1600 & molto scarso & $<$ 2000 & scarso \\
    \hline
    1600--2000 & scarso & $<$ 2400 & mediocre \\
    \hline
    2000--2400 & discreto & $<$ 2800 & buono \\
    \hline
    2400--2800 & buono & $<$ 3200 & molto buono \\
    \hline
    $>$ 2800 & ottimo & $>$ 3200 & eccellente \\
    \hline
  \end{tabularx}
\end{table}

\rubrica{Test di Gacon}
\begin{itemize}
  \item \textbf{sede}: pista di atletica; \textbf{obiettivo}: velocità massimale aerobica;
  \item \textbf{modalità}: percorrere distanze stabilite e progressivamente crescenti in 45$''$, con
    15$''$ di recupero --- si incrementa cioè la \textbf{distanza} a tempi fissi di fase attiva e di
    recupero;
  \item \textbf{progressione}: primo tratto di 100\,m per i sedentari e 125\,m per i soggetti
    allenati; i tratti successivi sono incrementati di 6,25\,m;
  \item \textbf{criterio di arresto}: quando il soggetto non copre più la distanza nei 45$''$
    stabiliti, la velocità del tratto precedente si considera quella massimale aerobica.
\end{itemize}

\rubrica{Test di Léger}
È il \textbf{primo test a navetta}, quello da cui sono derivati tutti gli altri.
\begin{itemize}
  \item \textbf{tipo}: test incrementale a navetta su 20\,m, con segnale sonoro cadenzato a
    intervalli stabiliti;
  \item \textbf{obiettivo}: velocità massimale aerobica e $FC_{max}$;
  \item \textbf{modalità}: percorrere i 20\,m nell'intervallo scandito dal segnale acustico;
    partenza da 8,5\,km/h con incrementi di 0,5\,km/h al minuto;
  \item \textbf{criterio di arresto}: il soggetto termina quando non riesce a coprire la distanza
    per due volte di seguito (test \textbf{a esaurimento});
  \item \textbf{stima} (con $V_{max}$ velocità massimale raggiunta):
    \begin{itemize}
      \item soggetti $>$ 18 anni: $VO_2max = -24{,}4 + 6 \cdot V_{max}$;
      \item soggetti tra 6 e 18 anni: $VO_2max = 31{,}025 - 3{,}238 \cdot V_{max} + 0{,}1536 \cdot
        V_{max} \cdot \textrm{età}$.
    \end{itemize}
\end{itemize}

\rubrica{Un'avvertenza sulle formule di stima}
Le formule di conversione sono ricavate da \textbf{gruppi ristretti} e perdono quindi il carattere
di individualità: risalire da metri percorsi o velocità massimale aerobica a un $VO_2max$ in
ml/kg/min è complicato e il risultato non corrisponde mai davvero alla realtà. Il suggerimento è di
tenere in considerazione i \textbf{soli dati misurati}: la \textbf{velocità massimale aerobica}
registrata prima della fine del test e la \textbf{frequenza cardiaca massima}.

\rubrica{La frequenza cardiaca massima}
La misurazione della $FC_{max}$ in questi test è di importanza fondamentale per la valutazione del
\textbf{carico interno} dell'allenamento: la frequenza cardiaca \textbf{relativa}, espressa in
percentuale della massima, dà l'intensità dello sforzo sostenuto nelle varie fasi
dell'allenamento e della partita. L'argomento è ripreso nella parte del corso dedicata al calcio.

\rubrica{Yo-yo endurance test}
Test successivi al Léger, nati con la \textbf{scuola danese} (\textbf{Bangsbo}, 1997). Le versioni
sono \textbf{tre} --- \textit{endurance}, \textit{intermittent endurance} e \textit{intermittent
recovery} --- e il protocollo completo è ripreso nel cap.~\ref{cap:metabolismoaerobico}:
\begin{itemize}
  \item \textbf{tipo}: test incrementale a navetta su 20\,m, con segnale sonoro cadenzato;
  \item \textbf{obiettivo}: velocità massimale aerobica, $FC_{max}$ e numero di frazioni di navetta
    (2 $\times$ 20\,m) percorse;
  \item \textbf{modalità}: come il Léger, ma con partenza da 8\,km/h per i principianti e da
    11,5\,km/h per gli atleti esperti.
\end{itemize}

\rubrica{Yo-yo intermittent recovery test}
\begin{itemize}
  \item \textbf{tipo}: variante del precedente, con struttura 20\,m $\times$ 2 più \textbf{10$''$ di
    recupero attivo} in \textit{jogging} attorno a un terzo delimitatore posto a \textbf{5\,m}
    dietro la linea di partenza (quindi 5 $+$ 5\,m);
  \item \textbf{obiettivo}: gli stessi del precedente;
  \item \textbf{modalità}: partenza da 10\,km/h per i principianti e da 13\,km/h per gli atleti
    esperti.
\end{itemize}

% ---------------------------------------------------------------------------
\section{La soglia anaerobica}

\rubrica{Il concetto}
\begin{itemize}
  \item \textbf{Wassermann e McIlroy} (\textit{Am J Cardiology}, 1964) introducono il termine
    \textbf{soglia anaerobica} per definire il carico di lavoro massimo mantenibile utilizzando il
    \textbf{solo} metabolismo aerobico --- la potenza muscolare, o la velocità di corsa, sostenibile
    senza accumulo di acido lattico e quindi, teoricamente, all'infinito:
    \begin{itemize}
      \item ``anaerobico'' fa riferimento alla \textbf{carenza relativa di ossigeno} a livello
        muscolare;
      \item ``soglia'' indica il \textbf{punto esatto}, espresso come potenza, in cui il fenomeno
        accade: il punto limite oltre il quale il soggetto comincia ad accumulare acido lattico e,
        mantenendo quell'intensità, finisce per fermarsi;
    \end{itemize}
  \item \textbf{Mader} (1976) dimostra che il carico di lavoro più alto al quale non avviene
    accumulo di lattato coincide approssimativamente con quello in cui, durante un test
    incrementale, si registra una concentrazione di lattato ematico di \textbf{4\,mM/l}.
\end{itemize}

\rubrica{Il valore predittivo}
\begin{itemize}
  \item la soglia anaerobica è il parametro metabolico \textbf{più correlato con la prestazione
    nelle gare di lunga durata}, più del $VO_2max$, del peso corporeo o della percentuale di fibre
    muscolari lente;
  \item correlazioni positive tra velocità di soglia e velocità di gara:
    \begin{itemize}
      \item \textbf{maratona}: $r$ = 0,98 (Farrel et al., \textit{MSSE}, 1979), su 13 fondisti;
      \item \textbf{10\,km}: $r$ = 0,94 (Powers S.K. et al., \textit{RQES}, 1983);
    \end{itemize}
  \item il metabolismo aerobico è anche ciò che permette di \textbf{recuperare}, riassorbendo a
    media intensità l'acido lattico prodotto.
\end{itemize}

\rubrica{Il test di Conconi (1976)}
\begin{itemize}
  \item \textbf{principio}: il test si basa sulla \textbf{perdita di correlazione lineare} tra
    frequenza cardiaca e carico di lavoro (velocità di corsa);
  \item \textbf{lettura}: il \textbf{punto di flessione} della curva della FC corrisponde alla
    soglia anaerobica, con relativa determinazione della \textbf{velocità di soglia};
  \item \textbf{protocollo originale}: aumenti di velocità dopo aver percorso una data distanza
    (1--2$''$ ogni 200\,m);
  \item \textbf{esempio}: nel grafico riportato la soglia cade a 12,0\,km/h e 170\,bpm, sotto la
    $FC_{max}$ e sotto il punto di massima attivazione del meccanismo aerobico;
  \item \textbf{vantaggio}: bastano i \textbf{cardiofrequenzimetri}, che hanno reso il test alla
    portata di tutti e semplice da eseguire;
  \item \textbf{limite}: \textbf{non è sempre facile individuare il punto di flesso}. Il test resta
    importantissimo per le discipline di lunga durata come il ciclismo, molto meno per gli sport di
    squadra e per il calcio.
\end{itemize}

\rubrica{Il test di Mognoni}
\begin{itemize}
  \item \textbf{metodo}: un \textbf{singolo prelievo} di sangue capillare al termine di una corsa
    alla velocità di 13,5\,km$\cdot$h$^{-1}$ della durata di 6$'$; il valore di lattato ematico
    ottenuto si sostituisce nella $x$ di un'equazione che restituisce, soggetto per soggetto, la
    velocità corrispondente all'\textbf{OBLA} (\textit{onset of blood lactate accumulation}), cioè
    all'inizio dell'accumulo di lattato ematico;
  \item \textbf{campo di applicazione}: solo atleti con massima potenza aerobica superiore a quella
    dei sedentari ma inferiore a quella degli atleti di resistenza --- categoria in cui rientrano i
    \textbf{calciatori};
  \item \textbf{risultato}: la velocità di soglia è molto più strettamente correlata con la
    \textbf{lattacidemia} ($r$ = 0,91) che con la \textbf{frequenza cardiaca} ($r$ = 0,32) misurate
    al termine del test.
\end{itemize}

% ---------------------------------------------------------------------------
\section{Che cosa conta nel calciatore}

\begin{enumerate}
  \item l'\textbf{obiettivo primario} è saper ripetere \textbf{sforzi anaerobici lattacidi},
    sostenuti da unità motorie veloci e reiterati nel tempo, per tutti i 90 minuti;
  \item la soglia anaerobica può essere valutata anche nel giocatore, ma è un \textbf{obiettivo di
    importanza secondaria}: più che come parametro della fase attiva va letta come \textbf{velocità
    di recupero} --- tanto più alta, tanti più metri percorsi e tanto più rapido il recupero;
  \item di conseguenza le fasi di recupero vanno legate al \textbf{livello di soglia} più che al
    $VO_2max$;
  \item il $VO_2max$ del calciatore (50--55\,ml/kg/min) è su \textbf{valori medi} rispetto agli
    altri sport di squadra e non è realmente condizionante, a differenza degli atleti di lunga
    distanza, che raggiungono valori di 65--80\,ml/kg/min: in quegli sport, però, il fattore
    limitante della prestazione è tutt'altro.
\end{enumerate}
